% =====================================================================
%  chunks/fm.tex —— 卷首：关于底本（Lead 撰）
% =====================================================================
\chapter*{关于底本}
\addcontentsline{toc}{chapter}{关于底本}
\markboth{关于底本}{关于底本}

本译注本所依据的底本，是科学出版社「国外数学名著系列（影印版）」的一种（封面标注「（影印版）1」）：

\begin{center}
\small
S.~P.~Novikov (Ed.), \emph{Topology~I: General Survey},\\
Translated from the Russian by Boris Botvinnik and Robert Burns,\\
Encyclopaedia of Mathematical Sciences, Vol.~12,\\
Springer-Verlag, Berlin--Heidelberg--New York, 1996.\\
（科学出版社影印本，北京）
\end{center}

原书是俄文《Itogi nauki i tekhniki》（《科学与技术总结》）系列中的拓扑学综述卷，俄文版成书于 1986 年；英译自 1993 年开始，1995 年 11 月 Novikov 又补写了「英译本导言」与附录，英译本于 1996 年出版。据英译本著录，本书四章正文、结语、英译本导言与附录均由 Novikov 执笔（丛书则著录他为主编），故全书行文风格一致，前后呼应极密。

底本为影印扫描本，正文没有文字层，排印体例保留了原书特征。翻译时以下几点需要提请读者注意：

\begin{enumerate}[leftmargin=2.2em,itemsep=2pt,topsep=2pt]
\item \textbf{编号体例照抄原书。}原书的公式与陈述按「节号.序号」编号：第 2 章 §4 内的公式即 (4.1)、(4.2)……，第 3 章 §3 内即 (3.1)、(3.2)……；第 1 章没有小节，故其编号形如 (0.1)、(0.2)。例（Example）在原文中只给序数，译文中照抄，如「例 8」；同一节内重新起编时也照抄其重号。\textbf{同一编号在不同章会重复出现}（各章自 §1 起重新编号），引用时须连同章号一起读，这也是原书的引用习惯（「见 (3.1)」指本章 §3 的公式，只有写出「第 4 章 §5」才是跨章引用）。
\item \textbf{记号照原书。}原书本就用黑板体 $\R$、$\C$、$\Z$（如原书 p.7 的 $\R^3$、p.273 的 $\R^n$），本书照用；$\R P^n$、$\C P^n$ 原书写作 $RP^n$、$CP^n$，本书统一为 $\RP^n$、$\CP^n$。余者如 $S^n$、$D^n$、$\pi_n(X)$、$H_n(X)$、$H^n(X)$ 等均照原书。
\item \textbf{插图与图题。}插图一律从影印本原页裁切，图内原有的英文标注与符号保持原样，图号照原书作「图 2.16」这样（节号.序号），不另编号。原书图下另有题注文字的，照译；原书只印图号、题注已写进图内者，译本只补一句提示图意的短题，或径留空题，以便与原图逐图对照。全书表格亦照原书，不另编号（原书表格本无表号）。
\item \textbf{章末读书笔记为本译注本所加。}每章（含附录）末尾的「本章读书笔记」由译者撰写，综述该章内容、线索与前后关联，不属于原书文字；每节开头的方框「本节重点」亦为译注。两者都排入正文，以与原著正文相区别的版式呈现。
\item \textbf{译注以脚注排出。}凡需说明体例、补足背景、指出原书排印问题之处，均写成脚注（标注「译注」），随文排在该页页脚；凡指出「原书 p. $N$ 作 X，按上下文/数学应为 Y」，都给出判断依据，读者可对照原书核查。
\item \textbf{正文一律照原书排印，不作径改。}原书明显的排印错误（脱字、错下标、符号张冠李戴等）在正文中一律照抄原样，只在译注中说明其应有形式与理由；确需按上下文读通的地方，也在译注中交代。这样读者手中的译本与原书页次始终可以逐字对照。
\end{enumerate}

% [Tutor Note: 本书的引用习惯容易让初读者困惑：正文里说「见 (3.1)」指的是本章 §3 的公式，而不是第 3 章的公式；说「见第 4 章 §5」才是跨章引用。凡遇到只有括号数字而无章号的引用，一律按当前章的节号理解。]

% CHUNK fm DONE
